\subsection{Graph of an Absolute Value Function} 
We begin with the graph of the base absolute value function.
\begin{lpic}{}{8}
\setgraph{5}{5}{1}{5}
\setspace{1}
\begin{scope}[xshift=\value{cols}*2\linespace-6\linespace,yshift=-6\linespace]
\AxesLarge
\end{scope}
\regtext{1.5}{-1}{$x$};
\regtext{5.5}{-1}{$y=f(x)=\Abs{x}$};
\draw (0,-1) -- (8,-1);
\draw (3,0) -- (3,-8);
\foreach \x in {-3,-2,-1}
	\regtext{1.5}{-5-\x}{$\x$};
\foreach \x in {0,1,...,3}
	\regtext{1.5}{-5-\x}{\figdash$\x$};
\end{lpic}
Every absolute value function will have a similar ``V'' shape (or upside down ``V'' shape).
\begin{procedure}{Graphing an Absolute Value}
\begin{enumerate}[1.]
\item Determine the point of the ``V''---where the input of the absolute value is 0.
\item Find one point with integer coordinates to each side of the point.
\item Fill in the sides of the ``V'' shape.
\end{enumerate}
\end{procedure}
\begin{example}Graph $f(x)=-\frac{3}{2}\Abs{x-1}+4$.
\begin{lpic}{}{10}
\begin{scope}[xshift=\value{cols}*2\linespace-6\linespace,yshift=-6\linespace]
\setgraph{5}{5}{3}{5}
\setspace{1}
\AxesLarge
\end{scope}
\regtext{1.5}{-4}{$x$};
\regtext{5.5}{-4}{$y=f(x)$};
\draw (0,-4) -- (8,-4);
\draw (3,-3) -- (3,-7);
\end{lpic}
\end{example}